% !TeX TS-program = xelatex
% 虚构演示简历 — 仅用于 hellojobs 测试数据
\documentclass{resume-photo}
\usepackage{xeCJK}
\setCJKmainfont{Noto Serif CJK SC}
\usepackage{fontspec}
\usepackage{enumitem}
\setlist[itemize]{itemsep=0pt, topsep=1pt, parsep=0pt, partopsep=0pt}

\ResumeName{张伟}

\begin{document}

\ResumeContacts{
  138-0001-0001,%
  \ResumeUrl{mailto:zhangwei26@zju.edu.cn}{zhangwei26@zju.edu.cn},%
  \textnormal{浙江大学 | 软件工程 · 硕士 | 2023—2026}%
}

\ResumeTitle

\noindent 主攻 Java 与分布式微服务，熟悉 Spring Cloud、Kafka 与 MySQL 分库分表；在电商高并发场景完成接口 P99 从 420ms 降至 180ms 的优化实践。注重可观测性与文档沉淀，习惯在评审阶段拆解风险；参与线上复盘与跨团队联调，英语 CET-6，可阅读官方文档与 RFC（演示数据）。

\section{教育背景}
\ResumeItem
{浙江大学}
[\textnormal{软件学院 | 软件工程 | 硕士 | 排名前 10\%}]
[2023.09—2026.06(预计)]

\begin{itemize}
  \item 核心课程：分布式系统、中间件原理、高可用架构设计。
\end{itemize}

\ResumeItem
{浙江大学}
[\textnormal{计算机学院 | 软件工程 | 学士}]
[2019.09—2023.06]

\begin{itemize}
  \item 校级优秀毕业生，ACM 校赛银牌。
\end{itemize}


\ResumeItem
{实验室与助教}
[\textnormal{本科生科研 / 课程助教}]
[2022—2025]

\begin{itemize}
  \item 进入校级重点实验室参与课题，负责数据采集、清洗与可视化看板搭建。
  \item 担任《数据结构》《操作系统》课程助教，批改作业 200+ 份，答疑课时 40h+。
  \item 组织期中复习串讲与上机辅导，课程评教得分位于前 10\%（演示）。
  \item 整理课程 FAQ 与实验踩坑文档，被后续年级沿用。
\end{itemize}



\section{专业技能}
\begin{itemize}
  \item 熟悉 Java、JVM 调优与 Spring Boot / Spring Cloud Alibaba 微服务栈。
  \item 掌握 MySQL 索引与事务、Redis 缓存与分布式锁、RocketMQ 异步解耦。
  \item 了解 Kubernetes 部署与 SkyWalking 链路追踪，具备线上问题定位经验。
  \item 熟悉 Linux 日常运维与 Shell，具备日志检索与 CPU/内存突增初步定位能力。
  \item 掌握 Git / CI 与 Code Review，了解 Docker 与 Kubernetes 基本编排与滚动发布。
  \item 了解 OAuth2 / JWT、接口幂等与 REST 规范，具备单元测试与集成测试习惯（JUnit/pytest 等）。
  \item 能阅读英文技术文档，具备跨团队沟通与需求澄清能力（测试简历）。
\end{itemize}

\section{实习经历}
\ResumeItem
{阿里巴巴集团 | 本地生活}
[\textnormal{Java 开发实习生}]
[2024.07—2025.01]

\begin{itemize}
  \item \textbf{职责}: 参与营销中台订单接口改造，日均调用量约 2 亿次。
  \item \textbf{成果}: 通过批量查询与本地缓存，核心下单链路 P99 延迟下降 58\%。
\end{itemize}


\ResumeItem
{小米集团 | 互联网商业部}
[\textnormal{广告服务端实习生}]
[2023.12—2024.05]

\begin{itemize}
  \item \textbf{业务}: 参与广告召回与排序链路监控埋点，建设异常流量识别规则。
  \item \textbf{性能}: 优化 protobuf 序列化热点，CPU 占用下降约 8\%。
  \item \textbf{数据}: 与数据团队协作搭建小时级报表，支持运营策略迭代。
  \item \textbf{稳定性}: 大促前完成容量评估与限流预案演练 3 轮。
  \item \textbf{总结}: 输出「大促保障清单」模板并在组内复用。
\end{itemize}



\section{项目经历}
\ResumeItem
{秒杀系统课程设计}
[\textnormal{高并发库存与限流}]
[2024.03—2024.06]

\begin{itemize}
  \item 基于 Redis 预减库存 + Lua 脚本保证原子性，支撑压测峰值 1.2 万 QPS。
  \item 网关层令牌桶限流，误杀率控制在 0.3\% 以下。
\end{itemize}


\ResumeItem
{API 网关插件体系}
[\textnormal{鉴权 / 限流 / 审计（演示）}]
[2024.01—2024.06]

\begin{itemize}
  \item \textbf{扩展}: 定义插件 SPI，支持 Lua / WASM 沙箱（演示）。
  \item \textbf{安全}: 统一 JWT 校验与 mTLS 可选链路。
  \item \textbf{运维}: 插件热更新与版本回滚，0 停机。
  \item \textbf{指标}: 网关自身 CPU 开销 <3\% 额外基线（压测）。
\end{itemize}



\section{个人荣誉}
\begin{itemize}
  \item 中国高校计算机大赛 网络技术挑战赛 华东赛区 一等奖
  \item 浙江大学 学业优秀奖学金 二等
  \item 全国计算机等级考试四级（网络工程师）（演示）
  \item 校级「优秀学生干部 / 优秀团员」或技术社团骨干（综合测评前 15\%）
  \item 开源贡献：向知名项目提交被合并的 PR（文档/测试/feature）（演示）
\end{itemize}

\end{document}
